\section{Hamilton's principle and Hamiltonian mechanics}

Newton's laws are often written in Cartesian coordinates, but a physical system may be described more naturally by angles, distances along constraints, or collective coordinates. Such variables are called generalized coordinates $q_i(t)$; they need not have dimensions of length. For generalized coordinates and velocities $\dot q_i(t)$, the action is
\begin{equation}
S[q]=\int_{t_1}^{t_2}L(q_i,\dot q_i,t)\,\mathrm{d}t.
\end{equation}
Hamilton's principle states that, among nearby paths with the same configurations at $t_1$ and $t_2$, the physical path satisfies
\begin{equation}
\delta S=0.
\end{equation}
The variation is with respect to the whole history $q_i(t)$, not an instantaneous displacement. Applying the Euler--Lagrange equation to the action gives
\begin{equation}
\frac{\mathrm{d}}{\mathrm{d}t}\frac{\partial L}{\partial\dot q_i}
-\frac{\partial L}{\partial q_i}=0.
\end{equation}
For $L=\frac12m\dot{\mathbf r}^{\,2}-V(\mathbf r)$, kinetic energy minus potential energy, this becomes
\begin{equation}
m\ddot{\mathbf r}=-\nabla V=\mathbf F.
\end{equation}

The Lagrangian formulation is coordinate-flexible and makes symmetries particularly visible. A Legendre transformation replaces the velocity variables by canonical momenta, defining
\begin{equation}
p_i=\frac{\partial L}{\partial\dot q_i},
\qquad
H(q_i,p_i,t)=p_i\dot q_i-L.
\end{equation}
The first relation may agree with mechanical momentum in simple Cartesian systems, but canonical momentum is defined by the Lagrangian and can differ when velocity-dependent interactions or curvilinear coordinates are present. Provided the relations $p_i=\partial L/\partial\dot q_i$ can be inverted for the velocities, the Hamiltonian is a function of $(q_i,p_i,t)$ on phase space. Taking the differential and comparing coefficients of $
\mathrm{d}q_i$, $\mathrm{d}p_i$, and $\mathrm{d}t$ gives Hamilton's equations:
\begin{equation}
\dot q_i=\frac{\partial H}{\partial p_i},
\qquad
\dot p_i=-\frac{\partial H}{\partial q_i},
\qquad
\frac{\partial H}{\partial t}=-\frac{\partial L}{\partial t}.
\end{equation}
The two first-order Hamilton equations are equivalent to the second-order Euler--Lagrange equations, but they treat coordinates and momenta symmetrically. If $L$ has no explicit time dependence, then $H$ is conserved. For ordinary mechanical systems, $H$ is the total energy when the Legendre transform is nonsingular; in more general systems it remains the generator of time evolution but need not equal a simple kinetic-plus-potential energy.